\documentclass[a4,11pt]{aleph-notas}
% Se puede ver la documentación aquí:
% https://github.com/alephsub0/LaTeX_aleph-notas

% -- Paquetes adicionales
\usepackage{enumitem}
\usepackage{aleph-comandos}
\usepackage{booktabs}


% -- Datos
\institucion{Facultad de Ciencias Exactas, Naturales y Ambientales}
\carrera{Catálogo STEM}
\asignatura{Métodos Numéricos}
\tema{Resumen no. 1: Fundamentos y método de bisección}
\autor{Mario Cueva - Andrés Merino}
\fecha{Periodo 2026-1}

\logouno[0.18\textwidth]{Logos/logoPUCE_04_ac}
\definecolor{colortext}{HTML}{1957d1}
\definecolor{colordef}{HTML}{1957d1}
\fuente{montserrat}


% -- Comandos adicionales
\setlist[enumerate]{label=\roman*.}


\begin{document}

\encabezado


%%%%%%%%%%%%%%%%%%%%%%%%%%%%%%%%%%%%%%
\section{Cifras significativas}

\begin{defi}
    Las \textbf{cifras significativas} de un número son los dígitos que contribuyen a su precisión. Estas incluyen todos los dígitos distintos de cero, los ceros entre dígitos significativos y los ceros finales en números decimales.
\end{defi}

\begin{ejem}
    El número 0.004560 tiene cuatro cifras significativas: 4, 5, 6 y el cero final.
\end{ejem}

\begin{ejem}
    El número 1000 tiene una cifra significativa si se escribe sin notación científica, pero si se escribe como $1.000 \times 10^3$, entonces tiene cuatro cifras significativas.
\end{ejem}

%%%%%%%%%%%%%%%%%%%%%%%%%%%%%%%%%%%%%%
\section{Sistema numérico}

Históricamente, se han desarrollado diversos sistemas numéricos para representar cantidades. Han existido sistemas como el romano, que entra en una categoría de sistemas no posicionales, y sistemas como el decimal, binario y hexadecimal, que son sistemas posicionales.

La característica principal de un sistema posicional es que el valor de cada dígito depende de su posición en el número. Así, en el número 343, el dígito 3 en la posición de las centenas representa 300, mientras que el dígito 3 en la posición de las unidades representa solo 3. Esto lo podemos expresar matemáticamente como:
\[
343 = 3 \times 100 + 4 \times 10 + 3 \times 1 = 3 \times 10^2 + 4 \times 10^1 + 3 \times 10^0.
\]
En el sistema decimal, la base es 10 y el conjunto de símbolos que pertenece a este sistema es $D = \{0, 1, 2, 3, 4, 5, 6, 7, 8, 9\}$.

En caso de ocupara otra base, como el binario (base 2), el conjunto de símbolos sería $D = \{0, 1\}$ y el número 1011 en binario se puede convertir a decimal de la siguiente manera:
\[
1011_2 = 1 \times 2^3 + 0 \times 2^2 + 1 \times 2^1 + 1 \times 2^0 = 8 + 0 + 2 + 1 = 11_{10}.
\]

\begin{defi}[Sistema numérico]
    Dado una base $b$ y un conjunto de dígitos $D = \{0, 1, \ldots, b-1\}$, un número entero en el sistema numérico se representa como una secuencia de dígitos de $D$.
    \[
    N = d_n d_{n-1} \ldots d_1 d_0 \ldots = \sum_{i=0}^{n} d_i b^i
    \]

    Para números decimales, se pueden incluir dígitos después de un punto decimal, representando fracciones de la base:
    \[
    N = d_n d_{n-1} \ldots d_1 d_0 . d_{-1} d_{-2} \ldots = \sum_{i=-\infty}^{n} d_i b^i.
    \]
\end{defi}

% Representación de enteros en computadoras

En una computadora, se utiliza el sistema binario. Para enteros, se puede utilizar $n$ bits, donde el primer bit se reserva para el signo (0 para positivo y 1 para negativo) y los $n-1$ bits restantes se utilizan para representar el valor absoluto del número. Esto se conoce como representación de enteros con signo.

\begin{defi}[Representación de enteros con signo]
    En una representación de enteros con signo utilizando $n$ bits, el rango de números que se pueden representar es de $-2^{n-1}$ a $2^{n-1}-1$.
\end{defi}

\begin{advertencia}
    Para $n=16$, el rango de enteros con signo es de $-32\,768$ a $32\,767$. Para $n=32$, el rango es de $-2\,147\,483\,648$ a $214\,483\,647$. Para $n=64$, el rango es de $-9\,223\,372\,036\,854\,775\,808$ a $9\,223\,372\,036\,854\,775\,807$.
\end{advertencia}


%%%%%%%%%%%%%%%%%%%%%%%%%%%%%%%%%%%%%%
\section{Definición de errores}

\begin{defi}[Error]
    El \textbf{error} es la diferencia entre un valor verdadero o exacto y un valor aproximado. En métodos numéricos, los errores aparecen porque los cálculos suelen trabajar con aproximaciones y con una cantidad finita de dígitos.
\end{defi}

\begin{defi}[Error absoluto]
    Sea $x$ el valor verdadero y $\widehat{x}$ una aproximación de $x$. El \textbf{error absoluto} se define como
    \[
        E_a = |x-\widehat{x}|.
    \]
    Este error mide la distancia entre el valor exacto y el valor aproximado.
\end{defi}

\begin{defi}[Error relativo]
    Sea $x \neq 0$ el valor verdadero y $\widehat{x}$ una aproximación de $x$. El \textbf{error relativo} se define como
    \[
        E_r = \frac{|x-\widehat{x}|}{|x|}.
    \]
    Si se desea expresarlo como porcentaje, se utiliza
    \[
        E_r\% = \frac{|x-\widehat{x}|}{|x|}\times 100.
    \]
\end{defi}

\begin{ejem}
    Si el valor verdadero es $x=3.1416$ y la aproximación es $\widehat{x}=3.14$, entonces
    \[
        E_a = |3.1416-3.14| = 0.0016
    \]
    y
    \[
        E_r\% = \frac{0.0016}{3.1416}\times 100 \approx 0.0509\%.
    \]
\end{ejem}

\begin{defi}[Error aproximado]
    Cuando no se conoce el valor verdadero, se puede estimar el error comparando dos aproximaciones consecutivas. Si $x_i$ es la aproximación actual y $x_{i-1}$ la aproximación anterior, el \textbf{error aproximado} se define como
    \[
        E_{\text{aprox}} = |x_i-x_{i-1}|.
    \]
    También se puede expresar de forma relativa como
    \[
        E_{\text{aprox},r}\% = \frac{|x_i-x_{i-1}|}{|x_i|}\times 100,
    \]
    siempre que $x_i \neq 0$.
\end{defi}

\begin{advertencia}
    Un error absoluto pequeño no siempre significa que la aproximación sea buena. Por ejemplo, un error de $0.01$ puede ser pequeño si se aproxima una cantidad cercana a $1000$, pero puede ser grande si se aproxima una cantidad cercana a $0.02$.
\end{advertencia}

%%%%%%%%%%%%%%%%%%%%%%%%%%%%%%%%%%%%%%
\section{Resolución de ecuaciones: Método de la bisección}

\begin{defi}[Raíz de una ecuación]
    Dada una función $f$, una \textbf{raíz} o \textbf{cero} de la ecuación $f(x)=0$ es un valor $r$ tal que
    \[
        f(r)=0.
    \]
    Resolver una ecuación no lineal consiste en encontrar uno o varios valores de $r$ que satisfagan esta condición.
\end{defi}

\begin{teo}[Cambio de signo]
    Si $f$ es continua en el intervalo $[a,b]$ y se cumple que
    \[
        f(a)f(b)<0,
    \]
    entonces existe al menos una raíz de $f(x)=0$ en el intervalo $(a,b)$.
\end{teo}

\begin{defi}[Método de la bisección]
    El \textbf{método de la bisección} es un método iterativo para aproximar una raíz de una función continua. Parte de un intervalo $[a,b]$ donde existe un cambio de signo y lo divide en dos subintervalos mediante el punto medio
    \[
        m = \frac{a+b}{2}.
    \]
    Luego se conserva el subintervalo donde se mantiene el cambio de signo.
\end{defi}

\begin{advertencia}
    Para aplicar el método de la bisección se necesita que la función sea continua en $[a,b]$ y que $f(a)f(b)<0$. Si no hay cambio de signo, el método no garantiza la existencia de una raíz en el intervalo.
\end{advertencia}

\begin{defi}[Iteración del método de la bisección]
    En cada iteración se calcula el punto medio
    \[
        m_i = \frac{a_i+b_i}{2}.
    \]
    Si $f(a_i)f(m_i)<0$, la raíz se encuentra en $[a_i,m_i]$ y se toma
    \[
        a_{i+1}=a_i, \qquad b_{i+1}=m_i.
    \]
    Si $f(m_i)f(b_i)<0$, la raíz se encuentra en $[m_i,b_i]$ y se toma
    \[
        a_{i+1}=m_i, \qquad b_{i+1}=b_i.
    \]
    Si $f(m_i)=0$, entonces $m_i$ es una raíz exacta.
\end{defi}

\begin{defi}[Criterios de parada]
    El método se detiene cuando la aproximación cumple una tolerancia establecida. Algunos criterios comunes son:
    \begin{enumerate}
        \item $|f(m_i)|<\varepsilon$,
        \item $|m_i-m_{i-1}|<\varepsilon$,
        \item $\dfrac{b_i-a_i}{2}<\varepsilon$,
        \item alcanzar un número máximo de iteraciones.
    \end{enumerate}
    El valor $\varepsilon>0$ representa la tolerancia aceptada para el error.
\end{defi}

\begin{ejem}
    Consideremos la ecuación
    \[
        f(x)=x^2-2=0.
    \]
    En el intervalo $[1,2]$ se tiene
    \[
        f(1)=-1, \qquad f(2)=2.
    \]
    Como $f(1)f(2)<0$, existe una raíz en $(1,2)$. El primer punto medio es
    \[
        m_1=\frac{1+2}{2}=1.5.
    \]
    Como $f(1.5)=0.25$, se cumple que $f(1)f(1.5)<0$, por lo que la raíz se encuentra en $[1,1.5]$. El siguiente punto medio es
    \[
        m_2=\frac{1+1.5}{2}=1.25.
    \]
    Repitiendo este proceso, los puntos medios se acercan a la raíz $\sqrt{2}\approx 1.4142$.
\end{ejem}

\begin{advertencia}
    La bisección es un método seguro cuando se cumplen sus condiciones, pero puede ser lento porque reduce el intervalo a la mitad en cada iteración.
\end{advertencia}

\end{document}
